\documentclass[11pt]{article}

\usepackage[a4paper,margin=0.9in]{geometry}
\usepackage{amsmath,amssymb}
\usepackage{enumitem}
\usepackage{parskip}
\usepackage{fancyhdr}

\pagestyle{fancy}
\fancyhf{}
\lhead{DA3410}
\rhead{Viva Answer Key}
\cfoot{\thepage}

\setlength{\parindent}{0pt}
\setlist[enumerate]{leftmargin=1.6em,itemsep=6pt,topsep=4pt}

\title{\LARGE DA3410 Viva Question Bank --- Answer Key}
\author{}
\date{}

\begin{document}
\maketitle
\thispagestyle{fancy}

\section*{Set 1}
\begin{enumerate}

\item \textbf{mean() vs.\ sum() in the loss.}
Answer: the loss value grows by a factor of $N$ (the batch size), and the
gradients scale by $N$ too.
Reason: \texttt{sum()} adds all $N$ per-sample terms instead of averaging
them, so both loss and gradient magnitude scale with batch size.

\item \textbf{ReLU at $x=3$.}
Answer: $f(x)=3$, $f'(x)=1$.
Reason: for $x>0$, ReLU passes the input through unchanged and has slope $1$.

\item \textbf{can two affine layers collapse to one?}
Answer: yes.
Reason: with no nonlinearity between them, $y=(xW_1+b_1)W_2+b_2=xW+b$ where
$W=W_1W_2$ and $b=b_1W_2+b_2$. A composition of affine maps is affine.

\item \textbf{dropping the max-subtraction in softmax.}
Answer: \texttt{exp(1000)} and \texttt{exp(1001)} overflow to \texttt{inf}
in float32, so \texttt{p} becomes \texttt{nan} (\texttt{inf/inf}).
Reason: subtracting the max keeps exponent arguments $\le 0$; without it,
large logits overflow before normalization.

\item \textbf{checkpointing bug.}
Answer: \texttt{best\_loss = float("inf")} sits inside the epoch loop, so it
resets every epoch and the check always passes.
Reason: the running best must be initialized once, before the loop, not on
every iteration.

\item \textbf{backward loop order.}
Answer: correct as written.
Reason: iterating \texttt{reversed(ops)} visits square, then relu, then
multiply -- the correct reverse-topological order for backpropagation.

\item \textbf{shape of grad\_W.}
Answer: $(10,7)$.
Reason: \texttt{x.T} is $(10,16)$ and \texttt{grad\_z} is $(16,7)$; matrix
product gives $(10,7)$.

\item \textbf{print output, $N=300$, batch\_size$=300$.}
Answer: \texttt{1 300}.
Reason: \texttt{range(0,300,300)} yields one start (0), so one update, and
that single batch has all 300 rows.

\item \textbf{checkpoint keeps changing.}
Answer: \texttt{best\_params = params} stores a reference, not a copy.
Reason: since \texttt{params} keeps being updated in place, the saved
"best" params silently changes along with it. Fix: use
\texttt{clone\_params(params)}.

\item \textbf{does the group split leak?}
Answer: no, it is safe.
Reason: groups are shuffled and partitioned first, and samples are then
selected by group membership, so no group can land in both splits.

\end{enumerate}

\section*{Set 2}
\begin{enumerate}

\item \textbf{evaluation bug.}
Answer: the scaler is fit on all of \texttt{X} before the train/val split.
Reason: this leaks validation-set statistics into the scaler; it should be
fit on the training split only, after splitting.

\item \textbf{ReLU at $x=0$.}
Answer: $f(x)=0$, $f'(x)=0$.
Reason: ReLU maps $0\mapsto 0$, and the stated convention sets the
derivative at zero to $0$.

\item \textbf{can this network collapse to one affine layer?}
Answer: no.
Reason: the ReLU between the two linear maps is a genuine nonlinearity, so
the composition is not affine in general.

\item \textbf{checkpointing bug.}
Answer: the condition checks \texttt{train\_loss} instead of
\texttt{val\_loss}.
Reason: checkpointing on training loss tracks fit to the training set, not
generalization, so it can save an overfit model as "best."

\item \textbf{backward loop bug.}
Answer: \texttt{op.parents[0].grad = op.grad} is wrong.
Reason: it overwrites the parent's gradient instead of accumulating
(\texttt{+=}), and only updates \texttt{parents[0]}, ignoring any other
parent (e.g.\ \texttt{multiply} has two).

\item \textbf{bias gradient reduction.}
Answer: shape $(128,)$; summing over \texttt{dim=0} is correct.
Reason: the bias is broadcast over the batch dimension (size 32) during the
forward pass, so its gradient must be summed over that same dimension to
match the bias shape.

\item \textbf{batches for $N=10$, batch\_size$=4$.}
Answer: 3 batches; last batch size 2.
Reason: starts are $0,4,8$; the last slice \texttt{X[8:12]} only has 2 rows
since $N=10$.

\item \textbf{ranking update noise.}
Answer: SGD $>$ mini-batch $>$ full-batch.
Reason: smaller batches give noisier gradient estimates; full-batch uses
the entire dataset, giving the least noise.

\item \textbf{updating the last layer before backprop finishes.}
Answer: this is incorrect.
Reason: earlier layers' gradients get computed using post-update weights,
which are inconsistent with the loss that was actually evaluated. All
gradients must be computed from the same forward pass before any weights
are changed.

\item \textbf{does this split leak?}
Answer: no, it is safe.
Reason: \texttt{group\_id != 5} and \texttt{group\_id == 5} are
complementary, so no sample can appear in both sets.

\end{enumerate}

\section*{Set 3}
\begin{enumerate}

\item \textbf{what data fits the scaler/PCA?}
Answer: training data only.
Reason: fitting on validation data leaks its statistics into
preprocessing, giving an overly optimistic evaluation.

\item \textbf{Leaky ReLU, $x=-2$, $\alpha=0.1$.}
Answer: $f(x)=-0.2$, $f'(x)=0.1$.
Reason: for $x<0$, Leaky ReLU returns $\alpha x$ and has slope $\alpha$.

\item \textbf{sigmoid bug.}
Answer: \texttt{1 / 1 + torch.exp(-x)} computes $1 + e^{-x}$, not
$1/(1+e^{-x})$.
Reason: division binds tighter than addition, so parentheses are needed:
\texttt{1 / (1 + torch.exp(-x))}.

\item \textbf{validation loop bug.}
Answer: the loop calls \texttt{backward(loss)} and \texttt{update(params,
lr)}.
Reason: this updates the model on validation data, so the reported
validation loss no longer reflects held-out performance.

\item \textbf{is \texttt{saved = \{"a": a\}} enough for $y=a{\cdot}b$?}
Answer: no.
Reason: the local gradients are $\partial y/\partial a=b$ and
$\partial y/\partial b=a$; \texttt{b} was not saved, so it must also be
retained.

\item \textbf{$a=3,b=2,g=1$.}
Answer: $\partial L/\partial a = 2$, $\partial L/\partial b = 3$.
Reason: for $y=ab$, $\partial L/\partial a = g\cdot b$ and
$\partial L/\partial b = g\cdot a$.

\item \textbf{correct gradient for $z=x@W$.}
Answer: Candidate A, \texttt{grad\_x = grad\_z @ W.T}.
Reason: with $x:(N,D)$, $W:(D,H)$, $\texttt{grad\_z}:(N,H)$, only
$(N,H)\!\times\!(H,D)=(N,D)$ matches the required shape of \texttt{grad\_x}.

\item \textbf{batches for $N=8$, batch\_size$=4$.}
Answer: 2 batches; last batch size 4.
Reason: starts are $0,4$, and $N=8$ divides evenly, so both batches are
full.

\item \textbf{checkpoint keeps changing.}
Answer: \texttt{best\_params = params} stores a reference, not a copy.
Reason: same as Set 1 Q29 -- use \texttt{clone\_params(params)}.

\item \textbf{He init with ReLU.}
Answer: stable.
Reason: He initialization is derived to preserve activation/gradient
variance specifically for ReLU networks.

\end{enumerate}

\section*{Set 4}
\begin{enumerate}

\item \textbf{what data fits the scaler/PCA?}
Answer: training data only. Reason: same as Set 3 Q4 -- fitting on
validation data leaks information.

\item \textbf{softmax of equal logits $[0,0]$.}
Answer: $p=[0.5,\,0.5]$.
Reason: equal logits get equal exponentials, so softmax splits probability
evenly.

\item \textbf{dropping the max-subtraction in softmax.}
Answer: same overflow as Set 1 Q12 -- \texttt{exp(1000)},
\texttt{exp(1001)} become \texttt{inf}, giving \texttt{nan}.
Reason: no upper bound on the exponent argument without subtracting the
max.

\item \textbf{shape of pred; effect of \texttt{dim=0}.}
Answer: \texttt{pred} has shape $(8,)$. With \texttt{dim=0}, shape becomes
$(7,)$.
Reason: \texttt{dim=1} picks the top class per sample (over 7 classes);
\texttt{dim=0} instead picks, per class column, which of the 8 samples
scores highest -- not a class prediction anymore.

\item \textbf{validation loop bug.}
Answer: \texttt{val\_loss} is set equal to \texttt{train\_loss} instead of
using the loss computed in the validation loop.
Reason: the computed \texttt{loss} is discarded; the reported
\texttt{val\_loss} does not reflect validation performance at all.

\item \textbf{is \texttt{saved = \{"mask": z>0\}} enough for ReLU?}
Answer: yes.
Reason: the ReLU backward pass only needs to know where the input was
positive; \texttt{grad\_z = grad\_h * mask} is exactly that.

\item \textbf{$a=4,b=5,g=2$.}
Answer: $\partial L/\partial a = 10$, $\partial L/\partial b = 8$.
Reason: $\partial L/\partial a = g\cdot b = 2\cdot5$;
$\partial L/\partial b = g\cdot a = 2\cdot4$.

\item \textbf{print output, $N=128$, batch\_size$=1$.}
Answer: \texttt{128 1}.
Reason: \texttt{range(0,128,1)} gives 128 starts, so 128 updates, and every
batch (including the last) has exactly 1 row.

\item \textbf{batches for $N=9$, batch\_size$=8$.}
Answer: 2 batches; last batch size 1.
Reason: starts are $0,8$; \texttt{X[8:16]} only has 1 row since $N=9$.

\item \textbf{Xavier init with sigmoid.}
Answer: stable.
Reason: Xavier (Glorot) initialization is derived to preserve variance for
symmetric activations like sigmoid/tanh.

\end{enumerate}

\section*{Set 5}
\begin{enumerate}

\item \textbf{mean() vs.\ sum() in the loss.}
Answer: same as Set 1 Q7 -- loss and gradients both scale by the batch
size $N$.

\item \textbf{softmax of equal logits $[-3,-3,-3,-3]$.}
Answer: $p=[0.25,0.25,0.25,0.25]$.
Reason: equal logits always give a uniform distribution over the classes,
regardless of their common value.

\item \textbf{shape of pred; effect of \texttt{dim=0}.}
Answer: same as Set 4 Q13 -- shape $(8,)$, and \texttt{dim=0} would pick
the top sample per class instead of the top class per sample.

\item \textbf{validation loop bug.}
Answer: \texttt{val\_loss} is computed on \texttt{train\_loader}, not a
held-out set.
Reason: evaluating on the training data cannot measure generalization; a
separate validation loader is required.

\item \textbf{is \texttt{saved = \{"s": s\}} enough for
\texttt{sigmoid}?}
Answer: yes.
Reason: the sigmoid derivative can be written purely in terms of its own
output, $s(1-s)$, so saving $s$ is sufficient.

\item \textbf{$a=-2,b=3,g=1$.}
Answer: $\partial L/\partial a = 3$, $\partial L/\partial b = -2$.
Reason: $\partial L/\partial a = g\cdot b$, $\partial L/\partial b = g\cdot
a$.

\item \textbf{shape of grad\_W.}
Answer: same as Set 1 Q22 -- $(10,7)$.

\item \textbf{print output, $N=256$, batch\_size$=128$.}
Answer: \texttt{2 128}.
Reason: \texttt{range(0,256,128)} gives 2 starts, and $256$ divides evenly,
so the last batch is also full size 128.

\item \textbf{permutation computed but not used.}
Answer: \texttt{xb}/\texttt{yb} are sliced with \texttt{i} directly instead
of \texttt{perm[i:i+batch\_size]}.
Reason: \texttt{perm} is never applied to index \texttt{X} and \texttt{y},
so the data stay in their original order.

\item \textbf{$2\times$He init with ReLU.}
Answer: exploding.
Reason: doubling the He-scaled weights increases per-layer variance beyond
the level He init is designed for, so activations/gradients grow layer to
layer.

\end{enumerate}

\section*{Set 6}
\begin{enumerate}

\item \textbf{evaluation bug.}
Answer: same as Set 2 Q6 -- the scaler is fit before the split, leaking
validation statistics.

\item \textbf{ReLU at $x=-2$.}
Answer: $f(x)=0$, $f'(x)=0$.
Reason: for $x<0$, ReLU outputs $0$ and has slope $0$.

\item \textbf{softmax of equal logits $[-2,-2,-2]$.}
Answer: $p=[1/3,\,1/3,\,1/3]$.
Reason: equal logits give a uniform distribution, here over 3 classes.

\item \textbf{checkpointing bug.}
Answer: \texttt{best\_params = params} stores a reference, not a copy.
Reason: initialization is correctly outside the loop this time, but the
save itself still needs \texttt{clone\_params(params)}.

\item \textbf{shape of z, $B=4,D=3,H=5$.}
Answer: $(4,5)$.
Reason: \texttt{x @ W} is $(4,3)\times(3,5)=(4,5)$; adding bias $(5,)$
broadcasts without changing the shape.

\item \textbf{backward loop bug.}
Answer: iterating \texttt{ops} in forward order instead of
\texttt{reversed(ops)}.
Reason: backprop needs reverse-topological order (loss back to inputs) so
each op's upstream gradient is ready before it is used; forward order
breaks this.

\item \textbf{backward accumulation bug.}
Answer: the second assignment overwrites the first instead of adding to
it.
Reason: \texttt{x} feeds \texttt{L} through two paths, so by the
multivariable chain rule the gradients must be summed:
\texttt{grad\_x = 1.0 + 2*x}, not reassigned.

\item \textbf{print output, $N=300$, batch\_size$=128$.}
Answer: \texttt{3 44}.
Reason: starts are $0,128,256$; the last slice \texttt{X[256:300]} has 44
rows.

\item \textbf{ranking update noise.}
Answer: same as Set 2 Q28 -- SGD $>$ mini-batch $>$ full-batch.

\item \textbf{does this split leak?}
Answer: yes, it can.
Reason: the split is done at the sample level via a random permutation,
not at the group level, so samples sharing a \texttt{group\_id} can end up
on both sides. Groups must be split first, then samples assigned by group.

\end{enumerate}

\end{document}
